\section{Exchange pictures and pinned returns}\label{sec:pins}

We express the magnetization through colored exchange cycles, then turn
pinning probabilities into one-particle return probabilities.  Two changes
of law are involved.  Pinning removes the color factor from cycles that
meet the pins.  Exposing the cycles that avoid the pins leaves a conditional
exchange process on the remaining slots.  The last part of this section
identifies its first return moment with an unconditioned Markov walk in the
exposed environment.  These constructions are finite-dimensional; the
spatial estimates begin in Section~\ref{sec:geometry}.

\subsection{Spin slots and the exchange law}

Fix $\beta>0$ and a cube $\Lambda=\Lambda_N$.  Set $\ell=2S$ and
\[
 b_0=\max\{|z|_\infty:J(z)>0\},\qquad
 \partial\Lambda=\{x\in\Lambda:
       \dist_\infty(x,\mathbb Z^3\setminus\Lambda)\le b_0\},\qquad
 \Lambda^\circ=\Lambda\setminus\partial\Lambda.
\]
No interaction bond joins $\Lambda^\circ$ to the exterior of $\Lambda$.
The \emph{slots} are the elements of
$V_\Lambda=\Lambda\times\{1,\ldots,\ell\}$.  Above each unordered
bond $\{x,y\}\subset\Lambda$ with $J(y-x)>0$, include all $\ell^2$
unordered slot pairs.  Write $\mathcal E_\Lambda$ for the resulting
slot edges and set
\[
 J_e=J(y-x),\qquad c_e=\frac{J_e}{2}
       \quad(e\text{ above }\{x,y\}),\qquad
 w_\Lambda=\sum_{e\in\mathcal E_\Lambda}J_e.
\]

Under the \emph{base law} $\mathbb P_{\rm b}$, every slot edge $e$
carries an independent Poisson process of rate $c_e$ on $[0,\beta)$.
At a mark the two slot lines exchange.  At the end of the interval,
independent uniform permutations of the slots at each site join time
$\beta$ to time $0$.  We call these within-site permutations the
\emph{seam completions}.  A picture $\omega$ consists of its marks and
completions.  Following its lines through repeated copies of the same
picture gives directed cycles.  Their intersections with the time-zero
layer are the cycles of a permutation of $V_\Lambda$.

Time belongs to $\mathbb R/\beta\mathbb Z$.  At an exchange we use the
slot just after that exchange, and at time zero we use the slot just
after the seam completion.  Deterministic times almost surely avoid
Poisson marks.  If $a_C$ is the number of time-zero slots on a cycle
$C$, then $\sum_Ca_C=\ell|\Lambda|$.  A cycle coloring is a choice
$\sigma_C\in\{-1,1\}$ for each cycle; its magnetization is
\[
                 M(\sigma)=\frac12\sum_Ca_C\sigma_C.
\]

\begin{proposition}[Loop representation]\label{prop:loop-representation}
For every finite cube, every $\beta>0$, and every real $h$,
\begin{equation}\label{eq:loop-partition}
 Z_\Lambda(h):=\Tr e^{-\beta(H^J_{S,N}-hM_{S,N})}
   =e^{\beta w_\Lambda/4}\,
           \mathbb E_{\rm b}\sum_\sigma e^{\beta hM(\sigma)}.
\end{equation}
At zero field the Gibbs spectral law of total magnetization is obtained
by weighting each picture by $2^{\#\{\text{cycles}\}}$, coloring its
cycles independently and fairly, and evaluating $M(\sigma)$.  More
precisely, for every function $\Phi$ on the finite spectrum of
$M_{S,N}$,
\begin{equation}\label{eq:loop-spectral-law}
 \frac{\Tr\bigl(e^{-\beta H^J_{S,N}}\Phi(M_{S,N})\bigr)}{Z_\Lambda(0)}
   =\frac{\mathbb E_{\rm b}\sum_\sigma\Phi(M(\sigma))}
           {\mathbb E_{\rm b}2^{\#\{\text{cycles}\}}}.
\end{equation}
\end{proposition}

\begin{proof}
This is the exchange representation of T\'oth~\cite{Toth1993}, with
the higher-spin symmetric-slot construction of
Nachtergaele~\cite[Section~2.2]{Nachtergaele1994}; see also
\cite[Lemma~2.1]{Ordering}.  At one site let $\psi_k$ be the normalized
sum of the elementary spin-$\tfrac12$ tensors having $k$ up spins,
$0\le k\le\ell$.  The sum of the slot $z$-spins acts on $\psi_k$
by $k-\ell/2$, and the sum of the raising operators acts by
$\sqrt{(\ell-k)(k+1)}\,\psi_{k+1}$.  These are the spin-$S$ matrices
under $k=S+m$.  Thus the physical space is the range of
\[
 \mathcal S_\Lambda=\prod_{x\in\Lambda}\mathcal S_x,
 \qquad \mathcal S_x=\frac1{\ell!}
                  \sum_{\kappa\in\mathfrak S_\ell}\mathcal U_{\kappa,x},
\]
where $\mathcal U$ denotes tensor permutation on the slot space.

The spin-$\tfrac12$ identity
$\boldsymbol s_i\cdot\boldsymbol s_j
 =\frac12\mathcal U_{(ij)}-\frac14I$ gives the lifted Hamiltonian
\[
 \widehat H=-\frac12\sum_eJ_e\mathcal U_e
                      +\frac14w_\Lambda I,
 \qquad \widehat M=\sum_{i\in V_\Lambda}s_i^z.
\]
The full sum over all slot pairs above each bond commutes with every
within-site permutation, and hence with $\mathcal S_\Lambda$.
The Poisson product $\varphi$ of the exchanges satisfies
\[
 \mathbb E_{\rm b}\mathcal U_\varphi
   =\exp\left\{\beta\sum_ec_e(\mathcal U_e-I)\right\}
   =e^{-\beta w_\Lambda/4}e^{-\beta\widehat H}.
\]
For example, the first equality follows from the matrix evolution with
generator $\sum_ec_e(\mathcal U_e-I)$; no commutation of individual
transpositions is needed.  Averaging the seam completions inserts
$\mathcal S_\Lambda$.  Taking the trace after multiplication by
$e^{\beta h\widehat M}$ proves \eqref{eq:loop-partition}: in the
elementary up/down basis, a tensor permutation has a nonzero diagonal
entry exactly when the coloring is constant on each cycle.  The same
calculation with any function of $\widehat M$ proves the spectral-law
assertion.
\end{proof}

\subsection{Pins, boundary approximation, and mark counts}

A \emph{pin} is a specified slot-time point.  Define the continuous
boundary pin set
\[
       D_0=\partial\Lambda\times\{1,\ldots,\ell\}
                                   \times(\mathbb R/\beta\mathbb Z).
\]
An \emph{allowed} pin set is $D=D_0\cup F$ with $F$ finite and
deterministic.  We also use finite pin sets before taking the boundary
limit.  For either type of set, put
\[
 n_D(\omega)=\#\{C:C\cap D=\varnothing\},\qquad
 \mathcal B_D(\omega)=\bigcup_{C:\,C\cap D=\varnothing}C,
 \qquad \mathcal Z_D=\mathbb E_{\rm b}2^{n_D},
\]
and define
\begin{equation}\label{eq:pinned-law}
 \mathbb P_D(d\omega)=\mathcal Z_D^{-1}2^{n_D(\omega)}
                                        \mathbb P_{\rm b}(d\omega).
\end{equation}
The normalizer obeys $1\le\mathcal Z_D\le2^{|V_\Lambda|}$.
In the associated zero-field color law, cycles meeting $D$ are up and
the remaining cycles receive independent fair colors.  Denote by
$d_D(i)$ the probability that a slot-time point $i$ is down.  Then
\begin{equation}\label{eq:down-free-cycle}
                  2d_D(i)=\mathbb P_D(i\in\mathcal B_D).
\end{equation}

\begin{lemma}[Dense boundary cuts]\label{lem:dense-boundary}
Let $F$ be a finite deterministic pin set and let $\mathcal T_n$ be
increasing finite sets of times containing zero, with dense union.
Set
\[
 D_n=F\cup\bigl(\partial\Lambda\times\{1,\ldots,\ell\}
                                      \times\mathcal T_n\bigr),
 \qquad D=F\cup D_0.
\]
Almost surely, $n_{D_n}=n_D$ and $\mathcal B_{D_n}=\mathcal B_D$ for
all sufficiently large $n$.  Moreover $\mathbb P_{D_n}$ converges to
$\mathbb P_D$ in total variation.  If bounded observables $X_n$
converge almost surely to $X$, with a common deterministic bound,
then $\mathbb E_{D_n}X_n\to\mathbb E_DX$.
\end{lemma}

\begin{proof}
There are finitely many marks per period, their times are distinct,
and seam completions preserve sites.  Every visit to the boundary
therefore contains a time interval of positive length and is detected
by a cut from the dense union.  There are finitely many cycles, so one
finite cut set eventually detects every boundary-touching cycle.
The asserted eventual equalities follow.  Since
$1\le2^{n_{D_n}}\le2^{|V_\Lambda|}$, bounded convergence applies both
to the weights and their positive normalizers.  This gives convergence
of the densities in $L^1(\mathbb P_{\rm b})$, and also the assertion
about observables.  This argument is at fixed volume.
\end{proof}

Under a pin law the marks need not be independent.  The following
uniform estimate on lists of distinct marks will suffice for all path
counts.  Let $\omega_{\rm m}$ be the mark counting measure on
$X_\Lambda=\mathcal E_\Lambda\times[0,\beta)$, and write
\[
             \lambda_{\rm b}(\{e\}\times dt)=c_e\,dt.
\]

\begin{lemma}[Factorial domination]\label{lem:factorial}
For an allowed or finite pin set $D$, an integer $j\ge1$, and a
nonnegative measurable function $F$ on $X_\Lambda^j$,
\begin{equation}\label{eq:factorial-bound}
 \mathbb E_D\sum_{(u_1,\ldots,u_j)\in(\omega_{\rm m})^j_{\ne}}
                         F(u_1,\ldots,u_j)
 \le 2^j\int F\,d\lambda_{\rm b}^{\otimes j}.
\end{equation}
Here the sum is over ordered tuples of distinct marks.  The estimate
also holds with any prescribed time order or time restrictions.
\end{lemma}

\begin{proof}
Inserting a transposition cuts and rejoins two strands.  It joins two
cycles or splits one cycle into two, leaving the other cycles
unchanged, and preserves the union of the affected strands and their
pin incidences.  Joining cycles cannot increase $n_D$.  On splitting,
a cycle avoiding $D$ produces two such cycles, while a cycle meeting
$D$ has at least one descendant that still meets $D$.  Thus $n_D$
increases by at most one, and the weight increases by at most two.
This statement
applies outside the null set of coincident marks or deterministic pin
times, and also to continuous boundary pins.

The Poisson insertion identity, with the seam fixed, is
\begin{align*}
 &\mathbb E_{\rm b}\left[2^{n_D(\omega)}
     \sum_{(u_1,\ldots,u_j)\in(\omega_{\rm m})^j_{\ne}}F(u_1,\ldots,u_j)
                                                        \right]\\
 &\qquad=\int F(u_1,\ldots,u_j)\,
       \mathbb E_{\rm b}2^{n_D(\omega+\sum_{a=1}^j\delta_{u_a})}
                          \prod_{a=1}^j\lambda_{\rm b}(du_a).
\end{align*}
Indeed, in the Poisson series the $k$-mark term has
$k!/(k-j)!$ ordered choices for the distinguished marks.  Cancelling
that factor against $1/k!$ and summing the remaining marks proves the
identity.  Nonnegativity justifies it even for infinite integrals.
The insertion bound and division by $\mathcal Z_D$ prove
\eqref{eq:factorial-bound}.
\end{proof}

We shall freely use the weaker constant $4^j$ in
\eqref{eq:factorial-bound}.  If an independent fresh Poisson family is
sampled as well, a mixed
list obeys the product of these bounds: each actual mark contributes
$4\lambda_{\rm b}$ and each fresh mark contributes
$\lambda_{\rm b}$.  This follows by independence and the two factorial
measures.  It is an unconditional statement under $\mathbb P_D$,
which is the law used when counting paths below.

\subsection{Single and simultaneous pin tests}

Let $E$ be allowed or finite, and let $T$ be a nonempty set of
$m_T=|T|$ distinct slot points at one deterministic time $\theta$,
disjoint from $E$.  Set $P=E\cup T$.  Repeating a sampled picture,
let $\mathcal R_i(E,T)$ be the event that the line from $i\in T$
visits $T$ at strictly positive elapsed time before visiting $E$.
Define
\begin{equation}\label{eq:pin-return-fraction}
 r(E,T)=\frac1{m_T}\sum_{i\in T}
                    \mathbb P_P\bigl(\mathcal R_i(E,T)\bigr).
\end{equation}
Pinning $T$ changes the picture law in this definition.  The next
proposition compares the changed law with the probability that $T$
lies on cycles avoiding the original pins.

\begin{proposition}[Pin identities]\label{prop:pin-identities}
For $E,T$ as above,
\begin{equation}\label{eq:pin-identities}
 d_E(i)=\frac{r(E,\{i\})}{1+r(E,\{i\})}\quad(i\notin E),
 \qquad
 \mathbb P_E(T\subseteq\mathcal B_E)\le(2r(E,T))^{m_T}.
\end{equation}
\end{proposition}

\begin{proof}
For $T=\{i\}$, write $r=r(E,\{i\})$ and
$A=\{i\in\mathcal B_E\}$. A line returns
to $i$ before $E$ exactly on $A$, so $r(E,\{i\})=\mathbb P_P(A)$.
Moreover $2^{n_E}=2^{n_P}(1+\ind{A})$, giving
$\mathcal Z_E=\mathcal Z_P(1+r)$ and
$\mathbb P_E(A)=2r/(1+r)$.  Equation~\eqref{eq:down-free-cycle}
proves the singleton identity.
This calculation applies directly to finite and allowed pin sets.

For the simultaneous bound, we use three precise finite-set facts from
\cite[Lemmas~3.2, 3.4, and 3.6]{Ordering}.  A nonzero polynomial is
\emph{stable} if it has no zero when every variable has positive
imaginary part.  The facts are as follows.
\begin{enumerate}[label=(\roman*),leftmargin=*]
\item Variable identification preserves stability.  Partial derivatives,
real specialization, and coefficientwise limits of uniformly bounded
degree preserve stability up to the zero polynomial.  In particular,
after identifying two variables as $t$, taking the coefficient of $t$
preserves stability up to zero.
\item If a permutation $\pi$ of a finite set $V$ is a product of
independent Bernoulli transpositions and deterministic permutations,
then
\[
            \mathbb E\prod_{j\in V}(x_j+y_{\pi(j)})
\]
is stable.  This also holds for independent products of finite-time
Poisson exchanges with arbitrary nonnegative rates and uniform
permutations on specified subsets.
\item If a random subset $A$ of a finite set has a stable generating
polynomial $\mathbb E\prod_{j\in A}y_j$, then for every subset $B$,
\[
                     \Pr(B\subseteq A)\le\prod_{j\in B}\Pr(j\in A).
\]
\end{enumerate}
These facts belong to the stable-polynomial method for negative
dependence developed in~\cite{BBL2009}.  They concern arbitrary finite
sets, with no spatial hypothesis.  To check the continuous-time scope
in (ii), subdivide a
finite interval into small slices and replace each edge clock in each
slice by its independent indicator of at least one mark.  Apply these
Bernoulli transpositions in a fixed order.  The products agree with the
Poisson products whenever no slice has two marks in total; the
probability of disagreement tends to zero.  A uniform permutation is
likewise a distributional limit of complete-graph Poisson exchanges:
the unit-time transition matrix on its finite permutation group has
all entries positive and preserves the uniform law, so it contracts
total variation distance to that law.  Thus the rates $c_e=J_e/2$
meet exactly the hypotheses of (ii).

For this bound, first suppose that $E$ is finite.  Cut the time circle at zero and
at every pin time in $P$.  Let $V$ be the disjoint union of the slot
layers at these cuts, and let $\pi$ follow a line from one cut to the
next cyclic cut.  Include the seam completion in the interval ending
at zero.  The interval maps are independent and use disjoint sets
of input $x$ variables and output $y$ variables.  Their stable symbols
therefore multiply to give the stable polynomial
\[
                  F_V=\mathbb E_{\rm b}
                         \prod_{j\in V}(x_j+y_{\pi(j)}).
\]

Delete an unpinned point $q\in V\setminus P$ by setting
$x_q=y_q=t$ and taking the linear coefficient in $t$.  For a fixed
permutation with predecessor $i$ and successor $j$ at a nonfixed $q$,
the two affected factors are $(x_i+t)(t+y_j)$, whose linear
coefficient is $x_i+y_j$.  This splices $i\to q\to j$ to $i\to j$.
A fixed point instead contributes the factor $2t$ and hence the
factor two.  Deleting all nonpins gives
\begin{equation}\label{eq:spliced-symbol}
 F_P=\mathbb E_{\rm b}\left[
         2^{n_P}\prod_{j\in P}(x_j+y_{\sigma(j)})\right],
\end{equation}
where $\sigma$ is the permutation taking each pin to its next pin
on the same cycle.  Each cycle avoiding $P$ contributes exactly one
factor two, when its final vertex is deleted.  The closure facts
give stability; the positive coefficients exclude the zero polynomial.

Differentiate \eqref{eq:spliced-symbol} in every $x_j$ with
$j\in P\setminus T$, set every remaining $x$ variable to zero, and
divide by $\mathcal Z_P$.  The result is the stable generating
polynomial
\begin{equation}\label{eq:pin-output-polynomial}
 G_T(\boldsymbol y)=\mathbb E_P\prod_{i\in T}y_{\sigma(i)}.
\end{equation}
It is nonzero because $G_T(\boldsymbol1)=1$, and it represents the
subset $\sigma(T)$ because its indices are distinct.  Counting the
outputs lying in $T$ gives
\[
 m_Tr(E,T)=\mathbb E_P|T\cap\sigma(T)|
                =\sum_{j\in T}\mathbb P_P(j\in\sigma(T)).
\]
Fact (iii) and the arithmetic--geometric mean inequality now imply
\[
              \mathbb P_P(T\subseteq\sigma(T))\le r(E,T)^{m_T}.
\]
The event $T\subseteq\sigma(T)$ is equivalent to
$T\subseteq\mathcal B_E$.  Indeed $|\sigma(T)|=|T|$, so the first
event says $\sigma(T)=T$; iterating $\sigma$ then meets only pins
of $T$ on every cycle through $T$, and meets all the pins on each
such cycle.  The converse follows by the same cycle description.

Deleting $m_T$ pins frees at most $m_T$ cycles.  Thus, pointwise,
\[
 2^{n_P}\le2^{n_E}\le2^{m_T}2^{n_P},
                   \qquad \mathcal Z_E\ge\mathcal Z_P>0.
\]
For $A=\{T\subseteq\mathcal B_E\}$, this yields
\[
 \mathbb P_E(A)
   \le 2^{m_T}\frac{\mathcal Z_P}{\mathcal Z_E}\mathbb P_P(A)
   \le (2r(E,T))^{m_T}.
\]

For $E=F\cup D_0$, replace $D_0$ by the finite cuts of
Lemma~\ref{lem:dense-boundary}, obtaining $E_n$ and $P_n=E_n\cup T$.
The simultaneous-free event and the weights converge almost surely.
The return events converge as well: follow each line from $i\in T$
until its first positive visit to $T$.  This horizon is finite because
the line eventually returns to $i$ on its finite permutation cycle.
Every boundary visit before that horizon contains a positive-duration
interval and is eventually detected.  Thus
$\ind{\mathcal R_i(E_n,T)}\to\ind{\mathcal R_i(E,T)}$.
Lemma~\ref{lem:dense-boundary} now passes the finite-pin inequality
to the limit.
\end{proof}

\subsection{Exposing the cycles that miss the pins}

The simultaneous pin inequality reduces the avoidance event to an average
return under $\mathbb P_P$.  We next separate the obstacles to that
return from the remaining exchanges.  From now on $E$ is allowed,
$T$ and $P$ are as above, and the outer picture law is $\mathbb P_P$.
Label a line by its slot at time zero.  Expose all labels belonging
to cycles missing $P$, their trajectories during one period, and their
endpoint mappings, including their seam images.  Write $\mathbf a$
for this exposure and $\mathscr H$ for its generated sigma-field.
At time $t$, the \emph{holes} $H_t=H_t(\mathbf a)$ are the slots
not occupied by exposed labels.  Their total number is constant.
Every pin is a hole at its pin time; every boundary slot is a hole
at every time.

For fixed $\mathbf a$, define a law $\mathbb Q_{\mathbf a}$ on the
remaining picture by the following operations.  While the hole set
is fixed, put independent rate-$c_e$ Poisson exchanges on the edges
whose endpoints are holes.  When an exposed label moves from a slot
to a hole, move that hole in the opposite direction into the vacated
slot.  Such a \emph{forced transport} is a fixed bijection between
the hole layers before and after the exposed jump.  Exchanges of two
exposed labels do not affect holes.  At the seam, independently at
each site, complete the revealed mapping by a uniform bijection
between its remaining source and target slots.  Let $G_P$ be the
event that every cycle of this remaining picture meets $P$, and put
\[
                         g_P(\mathbf a)=\mathbb Q_{\mathbf a}(G_P).
\]

\begin{proposition}[Exposure disintegration]\label{prop:exposure}
Almost surely under $\mathbb P_P$, $g_P(\mathbf a)>0$ and the
conditional remaining-picture law given $\mathscr H$ is
$\mathbb Q_{\mathbf a}(\,\cdot\mid G_P)$.

More precisely, for a deterministic initial label set
$A\subseteq V_\Lambda$, let $\mu_A$ be the base distribution of
its trajectories and endpoint mappings.  Let $F_A$ be the event
that these paths avoid $P$ and their endpoint set is $A$.  On $F_A$,
let $c_A(\mathbf a)$ be the number of cycles of their endpoint
permutation.  For any bounded measurable $\Phi$ of the exposure
and remaining picture $\omega_{\rm h}$,
\begin{equation}\label{eq:exposure-disintegration}
 \mathbb E_P\Phi(\mathbf a,\omega_{\rm h})
 =\frac1{\mathcal Z_P}\sum_{A\subseteq V_\Lambda}
     \int_{F_A}2^{c_A(\mathbf a)}
        \mathbb E_{\mathbb Q_{\mathbf a}}
                 [\ind{G_P}\Phi(\mathbf a,\omega_{\rm h})]
                                            \,\mu_A(d\mathbf a).
\end{equation}
In particular, the exposure marginal on this branch is
\begin{equation}\label{eq:exposure-marginal}
 \mathcal Z_P^{-1}\ind{F_A}2^{c_A(\mathbf a)}g_P(\mathbf a)
                                                   \,\mu_A(d\mathbf a).
\end{equation}
The same assertions hold for finite deterministic $P$.
\end{proposition}

\begin{proof}
We give the rate-dependent construction underlying
\cite[Proposition~4.1]{Ordering}.  Fix $A$ before sampling the base
picture.  On every edge $e$ take two independent rate-$c_e$ clocks.
Accept a mark from the first family when at least one endpoint
contains an $A$-label just before the mark; accept one from the second
family when neither endpoint does.

This produces the base exchange process.  Indeed, if
$c_{\rm tot}=\sum_ec_e>0$, the superposed clock has rate
$2c_{\rm tot}$.  At each of its times, conditionally on the past,
the accepted-edge or rejection outcome has probability
$c_e/(2c_{\rm tot})$ for edge $e$ and probability $1/2$ for
rejection.  These probabilities do not depend on the past or the
superposed waiting times.  Iterated conditioning and Poisson splitting
give independent accepted edge clocks of rates $c_e$.  If
$c_{\rm tot}=0$, the statement is immediate.

The $A$-label histories use only first-family clocks.  Given their
positions just before the seam, revealing their images under each
uniform site permutation leaves a uniform bijection of the remaining
slots, independently from site to site.  The histories, including
these revealed images, therefore do not use the second clock family
or the residual completions.  Conditional on those histories, the
second clocks restricted to the now fixed hole intervals remain
independent Poisson clocks, while the cross jumps give exactly the
forced transports.  This proves the base conditional kernel
$\mathbb Q_{\mathbf a}$ on $F_A$.

On $F_A$, the labels in $A$ form whole cycles missing $P$.  Requiring
$G_P$ makes them precisely all such cycles.  Consequently each
complete picture contributes to exactly one term in
\[
                    2^{n_P}
        =\sum_{A\subseteq V_\Lambda}\ind{F_A}2^{c_A}\ind{G_P},
\]
where a summand is zero off $F_A$.  Apply the base conditional kernel
to this finite sum and divide by $\mathcal Z_P$ to obtain
\eqref{eq:exposure-disintegration}.  Integrating out the remainder
gives \eqref{eq:exposure-marginal}, which gives zero outer mass to
$g_P=0$.  Dividing by $g_P$ on its positive set proves the asserted
conditional law.  All history spaces here are finite-jump path spaces
on finite state spaces, with finite endpoint data; they are standard
Borel spaces and admit the conditional kernels used above.
\end{proof}

The factor $g_P$ in \eqref{eq:exposure-marginal} is part of the
exposure law throughout the proof.  In particular, averaging a function
of the holes will always mean averaging under the actual $\mathbb P_P$
marginal.  At a fixed exposure, however, the next algebraic calculation
uses the unconditioned law $\mathbb Q_{\mathbf a}$.

\subsection{Removing the cycle condition from a first moment}

We record the finite-dimensional argument from
\cite[Lemma~4.2 and Proposition~4.3]{Ordering}, including its proof
to specify precisely why fresh transitions appear.  Matrices use the
column-forward convention: a permutation matrix obeys
$\Pi e_i=e_{\pi(i)}$.  For a matrix $M$, $\bigwedge^kM$ denotes
its action on alternating $k$-tensors; its matrix entries are the
$k\times k$ minors of $M$, and $\bigwedge^0M=1$.

\begin{lemma}[Exterior moments]\label{lem:exterior-moments}
Suppose a random permutation matrix $\Pi$ on a finite set is a
distributional limit of products of independent Bernoulli
transpositions and deterministic permutations.  If $K=\mathbb E\Pi$,
then
\begin{equation}\label{eq:exterior-moments}
            \mathbb E\bigwedge\nolimits^k\Pi=\bigwedge\nolimits^kK
                        \qquad(0\le k\le|V|).
\end{equation}
For every fixed exposure and finite set of cuts, the layered
permutation obtained under $\mathbb Q_{\mathbf a}$ by following
holes to the next cyclic cut has this property.
\end{lemma}

\begin{proof}
For a transposition matrix $T_{ij}$,
\[
                 T_{ij}-I=-(e_i-e_j)(e_i-e_j)^*.
\]
Every minor of $I+u(T_{ij}-I)$ is affine in $u$, since a term using
two columns of the rank-one perturbation has determinant zero.
Thus
\[
 \bigwedge\nolimits^k\bigl((1-u)I+uT_{ij}\bigr)
   =(1-u)\bigwedge\nolimits^kI+u\bigwedge\nolimits^kT_{ij},
\]
which proves the identity for one Bernoulli transposition.
Multiplicativity of exterior powers and independence preserve it
under products.  Deterministic permutations satisfy it directly,
and continuity of the minors on the finite matrix space passes it
to distributional limits.

For the hole picture, subdivide each interval between prescribed
changes and use the Bernoulli approximation already described.
Forced transports are deterministic bijections.  The residual seam
completions are uniform bijections, each expressible as a fixed
bijection composed with a uniform permutation; the latter has the
complete-graph exchange approximation described above.  Identify
every cut layer with a fixed set of $|H_t|$ elements.  The interval
maps act independently on separate copies of this set; their direct
sum followed by the deterministic cyclic shift of the layers is
the layered permutation.  A transposition in one block is a
transposition on the full layered set.  This verifies the required
product representation, before imposing $G_P$.
\end{proof}

\begin{proposition}[Conditional next-pin kernel]\label{prop:next-pin-kernel}
Let $V=I\sqcup B$ be a finite set with $I\ne\varnothing$.  Suppose
a random permutation matrix $\Pi$ satisfies
\eqref{eq:exterior-moments}, and let $K=\mathbb E\Pi$.  Let $G_I$
be the event that every cycle meets $I$.  On $G_I$, let $\Sigma$
be the column-forward permutation matrix on $I$ giving the next
positive visit to $I$.
Then
\begin{equation}\label{eq:all-hit-determinant}
                         \Pr(G_I)=\det(I-K_{BB}).
\end{equation}
If this probability is positive, then $K_{BB}$ is transient and
\begin{equation}\label{eq:conditional-next-pin}
 \mathbb E[\Sigma\mid G_I]
       =K_{II}+K_{IB}(I-K_{BB})^{-1}K_{BI}.
\end{equation}
The right-hand side is the first positive $I$-hitting kernel of
the Markov chain with independent transitions according to $K$.
If $B$ is empty, the determinant is one and the correction is zero.
\end{proposition}

\begin{proof}
The principal-minor expansion and \eqref{eq:exterior-moments} give,
for every fixed matrix $D$,
\begin{equation}\label{eq:determinant-expectation}
 \mathbb E\det(I-D\Pi)
  =\sum_k(-1)^k\Tr\left[(\bigwedge\nolimits^kD)
                           \mathbb E\bigwedge\nolimits^k\Pi\right]
  =\det(I-DK).
\end{equation}
Choose $D=\operatorname{diag}(W,I_B)$ with an arbitrary complex
$I\times I$ matrix $W$.  If a cycle lies in $B$, its indicator
vector is fixed by $D\Pi$, so the determinant vanishes.  On $G_I$,
$\Pi_{BB}$ is nilpotent and its resolvent follows the successive
nonpins on a permutation cycle.  Block elimination therefore gives
\[
 \det(I-D\Pi)
 =\det\left(I-W[\Pi_{II}+
               \Pi_{IB}(I-\Pi_{BB})^{-1}\Pi_{BI}]\right)
 =\det(I-W\Sigma).
\]
Consequently
\begin{equation}\label{eq:next-pin-determinant}
 \mathbb E[\ind{G_I}\det(I-W\Sigma)]
     =\det\begin{pmatrix}
             I-WK_{II}&-WK_{IB}\\
             -K_{BI}&I-K_{BB}
           \end{pmatrix}.
\end{equation}
Setting $W=0$ proves \eqref{eq:all-hit-determinant} without any
invertibility assumption.  When its value is positive,
$I-K_{BB}$ is invertible; taking the Schur complement in
\eqref{eq:next-pin-determinant} and dividing by that determinant yields
\[
 \mathbb E[\det(I-W\Sigma)\mid G_I]=\det(I-WH),
 \qquad H=K_{II}+K_{IB}(I-K_{BB})^{-1}K_{BI}.
\]
Replace $W$ by $tW$ and compare linear coefficients.  The identity
$\det(I-tWA)=1-t\Tr(WA)+O(t^2)$ shows that
$\Tr(W\mathbb E[\Sigma\mid G_I])=\Tr(WH)$ for every $W$.
Matrix units recover every entry, proving
\eqref{eq:conditional-next-pin}.

Finally, $K$ is nonnegative and doubly stochastic.  Its restriction
$K_{BB}$ is substochastic.  If some state of $B$ could never reach
$I$, its reachable set would contain a closed class and force the
eigenvalue one for $K_{BB}$, contradicting invertibility.  Finiteness
then supplies an integer $a$ and $\varepsilon>0$ such that every
state exits $B$ within $a$ steps with probability at least
$\varepsilon$.  Hence $\|K_{BB}^{an}\|_{1\to1}\le(1-\varepsilon)^n$
and
\[
                       (I-K_{BB})^{-1}=\sum_{n\ge0}K_{BB}^n.
\]
The first term in $H$ is a direct hit of $I$, and
$K_{IB}K_{BB}^nK_{BI}$ records a first hit after $n+1$ intervening
states of $B$.  This proves the asserted Markov interpretation.
\end{proof}

\subsection{The fresh walk}

Fix an exposure with $g_P(\mathbf a)>0$.  Repeat its hole sets,
forced transports, and revealed seam restrictions periodically.
The \emph{fresh walk} follows the one-particle transitions of
$\mathbb Q_{\mathbf a}$: it jumps across each available slot edge
at rate $c_e$, follows the forced transports, and at a seam uses a
fresh uniform residual completion.  Its random choices on disjoint
traversal intervals, including later repetitions of the same physical
interval, are independent.  Finite jump rates and finitely many
prescribed changes per period make this a nonexplosive Markov walk.
For $i\in T$, let $\tau_T^+$ be its first visit to $T$ at strictly
positive elapsed time and $\tau_E$ its first visit to $E$; an absent
visit has time infinity.  Put
\begin{equation}\label{eq:fresh-q}
             q_i(\mathbf a)=\Pr^{\rm walk}_{\mathbf a,i}
                                         (\tau_T^+<\tau_E).
\end{equation}

\begin{proposition}[Fresh-walk identity]\label{prop:fresh-walk}
For almost every exposure under $\mathbb P_P$ and every $i\in T$,
\[
       \mathbb P_P(\mathcal R_i(E,T)\mid\mathscr H)=q_i(\mathbf a).
\]
In particular,
\begin{equation}\label{eq:fresh-return}
                m_Tr(E,T)=\mathbb E_P\sum_{i\in T}q_i(\mathbf a).
\end{equation}
\end{proposition}

\begin{proof}
First use finitely many pins and take cuts at all pin times and
zero.  The unconditioned layered hole permutation under
$\mathbb Q_{\mathbf a}$ satisfies
Lemma~\ref{lem:exterior-moments}.  Conditional on every cycle hitting
the pins, Proposition~\ref{prop:next-pin-kernel} identifies its
next-pin first moment with independent transitions of its mean matrix.
Those transitions are exactly the fresh walk observed at the cuts.
Together with Proposition~\ref{prop:exposure}, this proves the
finite-pin assertion.

For an allowed $E=F\cup D_0$, fix an exposure for which the
disintegration holds and $g_P>0$.  Approximate the boundary by
increasing dense finite cuts and write $E_n=F\cup D_{0,n}$,
$P_n=E_n\cup T$.  The cuts include zero and all times in $P_n$.
Under the fixed unconditioned law $\mathbb Q_{\mathbf a}$, let
$G_n$ mean that every remaining cycle meets $P_n$, and put
$g_n=\mathbb Q_{\mathbf a}(G_n)$.

Every boundary visit of the remaining picture has positive residence
time.  Its Poisson marks almost surely avoid the finitely many
prescribed changes; seam completions preserve sites; and exposed
histories avoid the boundary, so forced transports cannot cause an
instantaneous boundary excursion.  Consequently
\[
              \ind{G_n}\uparrow\ind{G_P},\qquad
                         g_n\uparrow g_P>0.
\]
For all sufficiently large $n$, Proposition~\ref{prop:next-pin-kernel}
is therefore applicable with the pin layer $I=P_n$.  If $A_n$ is
the event that the next $P_n$-pin after $i$ is in $T$, it gives
\begin{equation}\label{eq:finite-cut-fresh-return}
 \frac{\mathbb Q_{\mathbf a}(G_n\cap A_n)}{g_n}
       =\Pr^{\rm walk}_{\mathbf a,i}(\tau_T^+<\tau_{E_n}).
\end{equation}
No assertion about the possibly singular earlier $n$ is needed.

In a fixed remaining picture the first positive return from $i$ to
$T$ is finite.  Detecting all boundary visits before that return
shows that $\ind{A_n}$ decreases to
$\ind{\mathcal R_i(E,T)}$.  Bounded convergence on the left of
\eqref{eq:finite-cut-fresh-return}, including its positive limiting
denominator, gives the return probability under
$\mathbb Q_{\mathbf a}(\,\cdot\mid G_P)$.

On the right use one uncensored fresh trajectory for all $n$.
The events decrease; every one requires $\tau_T^+<\infty$.
Any boundary visit before this finite time is detected by a dense
cut, again because it has positive duration.  Their intersection
is therefore exactly $\{\tau_T^+<\tau_E\}$ up to a null set.
Continuity from above proves convergence to \eqref{eq:fresh-q}.
The disintegration identifies the two limits, and averaging under
the exposure marginal \eqref{eq:exposure-marginal} proves
\eqref{eq:fresh-return}.
\end{proof}

We have obtained a Markov return problem in a periodic environment
whose missing slots are the exposed cycles.  The environment retains
its full pinned-law distribution, while the walk resamples its
transitions on every traversal.  The next section develops the
deterministic heat-flow estimates needed to bound these returns.
